\ifdefined\MyWebBook\else
\documentclass[../textbook.tex]{subfiles}
\begin{document}
\fi
\bookchapter{模型规模与资源预算}{ch:resource-budget}

一个模型能否在设备上运行，以及能否在期限内完成训练，是两个不同问题。前者取决于峰值状态和工作区，后者取决于总工作量与有效处理速度。只给参数量，既定不下显存，也定不下设备数。本章从模块形状出发，建立参数、计算、存储与期限之间可以逐项解释的关系。

资源预算以工作负载为起点，将参数量、计算量和状态存储换算为设备容量与时间需求。设备、网络和存储的供给条件见\bookxref{ch:compute-infrastructure}，状态切分与通信机制见\bookxref{ch:distributed-training}。

\begin{assumption}
除另行说明外，考虑自回归 Transformer，查询总宽度满足 $h_qd_h=d$，键和值等宽；线性层不含偏置，前馈层采用三个矩阵的门控结构。将一次乘法与一次加法计为两个浮点操作。参数存储、缓存、激活和工作区分别计量，所有内存先换算为字节后再求和。
\end{assumption}

\section{计量对象与单位}

\begin{table}[htbp]
\centering
\caption{资源与算力预算主要符号与计量对象}
\label{tab:rev-budget-symbols}
\begin{tabularx}{\linewidth}{lX}
\toprule 符号 & 含义\\\midrule
$L,d,V,d_f$ & 层数、隐藏宽度、词表大小、前馈中间宽度。\\
$h_q,h_{kv},d_h$ & 查询头数、键值头数、每头宽度。\\
$B,T,D$ & 一次计算的序列批量、序列长度、整个训练过程消费的词元数。\\
$P,F$ & 参数数与浮点操作数，二者量纲不同。\\
$b_w,b_g,b_c$ & 权重、梯度、缓存每元素字节数。\\
$N,S,u$ & 设备数、时间预算（秒）、按指定口径定义的有效利用率。\\
$R_{\mathrm{peak}},\beta$ & 单设备峰值操作率、有效数据传输带宽。\\
\bottomrule
\end{tabularx}
\end{table}

\term{浮点操作数}{Floating-point Operations}{FLOPs}表示完成给定工作量所需的操作数量；FLOP/s 才是每秒操作率。参数数是独立可学习标量的数量。共享同一矩阵两次使用不会增加独立参数，却可能增加计算次数。

十进制 GB 为 $10^9$ 字节，二进制 GiB 为 $2^{30}$ 字节。若模型含 $7\times10^9$ 个参数，每参数两字节，权重为 $14\times10^9$ 字节，约 $13.04$ GiB。把“7B”与“7 GiB”当成同一个数，会在预算第一步就产生错误。\footnote{硬件标称容量、操作系统显示和框架日志可能采用不同单位。有些界面虽写 GB，实际计算却用二进制除数，资源报告应列出原始字节数及换算规则。}

峰值还必须带时间语义。训练过程中权重、梯度、保存激活、通信缓冲和检查点暂存的生存区间不同，逐项峰值直接相加可能过于保守；只记录每项平均值又可能漏掉同时存在的瞬时峰值。初步预算可先给保守上界，再用执行时间线细化。

\section{参数量核算}

\subsection{注意力与门控前馈层}
查询投影 $\mat{W}_Q\in\Real^{d\times h_qd_h}$，键值投影 $\mat{W}_K,\mat{W}_V\in\Real^{d\times h_{kv}d_h}$，输出投影 $\mat{W}_O\in\Real^{h_qd_h\times d}$。因此
\begin{align}
 P_{\mathrm{attn}}
 &=d h_qd_h+2d h_{kv}d_h+h_qd_hd\\
 &=2d^2+2d h_{kv}d_h.\label{eq:budget-attn-params}
\end{align}
第二行使用了 $h_qd_h=d$；结构不满足这一条件时，第一行才是有效形式。GQA 减少键值投影参数，但查询和输出两项没有因此消失。

门控前馈采用两条上投影和一条下投影，矩阵形状分别为 $d\times d_f,d\times d_f,d_f\times d$，故
\begin{equation}
 P_{\mathrm{ffn}}=3dd_f.
\end{equation}
普通两矩阵前馈则为 $2dd_f$。比较时可固定中间宽度或参数预算，两种条件分别对应不同的中间维数和计算量。

若每层有两组只含缩放参数的 RMSNorm，末尾还有一组，则归一化参数为 $(2L+1)d$。若改为同时含缩放与偏置的 LayerNorm，需要按实际数量重新计数。偏置虽常是小项，精确统计仍应说明是否包含。

\subsection{输入输出共享与专家层}
输入嵌入参数为 $Vd$。输出投影若独立，再增加 $Vd$；若与输入共享，则独立存储只计一次。设 $\tau=1$ 表示共享，$\tau=2$ 表示独立，在上述同构稠密结构下
\begin{equation}\label{eq:budget-totalparams}
 P=\tau Vd+L(2d^2+2dh_{kv}d_h+3dd_f)+(2L+1)d.
\end{equation}
学习到的绝对位置表、额外任务头或查询键归一化参数不包含在该式中，存在时必须逐项加入。

对于专家层，若每个专家为门控前馈，独立路由专家总参数为 $E\cdot3dd_e$；每词元执行的专家矩阵规模近似为 $k\cdot3dd_e$。共享专家、路由器和注意力另计。\keyconcept{激活参数影响本次计算，总参数影响需要保存的模型状态。}显存公式中两者各自出现在不同位置。路由不均衡还会使某设备的实际负载偏离平均值，详见\bookxref{ch:rev-moe}。

\subsection{参数量核算}
\begin{example}


构造模型取 $L=12,d=1024,V=32000,h_q=16,h_{kv}=4,d_h=64,d_f=2816$，输入输出共享，归一化采用上述 RMSNorm 配置。注意力每层参数为
\begin{equation}
 2(1024)^2+2(1024)(4)(64)=2621440.
\end{equation}
前馈每层为 $3(1024)(2816)=8650752$。12层主干矩阵共 $135266304$ 个参数，嵌入为 $32768000$，归一化为 $25600$，合计
\begin{equation}
 P=168059904.
\end{equation}
若只把输出改为独立矩阵，增加 $32768000$ 个参数。若只把键值头由4改成16，每层增加 $2d(16-4)d_h$ 个参数；这一改动远小于把整个模型参数乘四。
\end{example}

\section{矩阵计算量与训练近似}

\subsection{线性层前向与反向}
对 $\mat{X}\in\Real^{n\times a}$ 和 $\mat{W}\in\Real^{a\times b}$，输出 $\mat{Y}=\mat{X}\mat{W}$ 有 $nb$ 个元素，每个元素执行 $a$ 次乘法和约 $a$ 次加法，故前向近似 $2nab$ FLOPs。该线性层有 $ab$ 个参数，因此每行输入的矩阵计算约为 $2P_W$。

反向需要 $\overline X=\overline YW^{\mathsf T}$ 与 $\overline W=X^{\mathsf T}\overline Y$，两项各约 $2nab$，前向与反向合计约 $6nab$。这解释了稠密模型训练中常用 $6PD$ 估计的来源：每词元约执行一次前向与两次同阶矩阵反向。各执行阶段的核心计算与显存速查规则如表\ref{tab:rev-compute-rules}所示。

\begin{table}[htbp]
\centering
\small
\setlength{\tabcolsep}{3.5pt}
\caption{自回归 Transformer 各执行阶段 FLOPs 与显存估算速查规则}
\label{tab:rev-compute-rules}
\begin{tabularx}{\linewidth}{lp{22mm}p{28mm}X}
\toprule
执行阶段 & 每步/每词元 FLOPs & 显存主要占用 & 关键假设与工程边界 \\
\midrule
前向传播 & $2P$ & 激活 $M_{\mathrm{act}} \propto L B T d$ & 密集矩阵乘法，不含 Softmax 与非线性 \\
反向传播 & $4P$（输入与权重梯度） & 梯度 $M_g = P \cdot b_g$ & 需反传更早层时，冻结层亦须求输入梯度 \\
稠密训练单步 & 约 $6PD$（$D$ 为词元数） & 权重、梯度与优化器 & 未计长序列 $\mathcal{O}(T^2)$ 项与激活重计算 \\
重计算训练 & 约 $8PD$（重算前向 $2P$） & 峰值降为 $\mathcal{O}(B T d)$ & 用约 $33\%$ 额外计算换取激活显存骤降 \\
自回归预填充 & $2P \cdot T$ & KV 缓存写入量 & 算力受限，吞吐由 Tensor Core 决定 \\
自回归解码 & $2P$（每生成一词元） & KV 缓存常驻与累加 & 访存受限，吞吐由显存带宽 $\beta$ 决定 \\
\bottomrule
\end{tabularx}
\end{table}

该估计针对参与矩阵乘法的参数。查表、共享参数的重复使用、冻结模块、优化器更新、专家路由和非矩阵算子需分别计量。尤其在门控网络中，逐元素乘法和激活导数有额外成本；大矩阵占主导时，把它们作为次阶项才有依据。

\subsection{注意力的长度平方项}
对每条序列，$QK^{\mathsf T}$ 的计算近似为 $2h_qT^2d_h$，再与值矩阵相乘也约为同量级，因此
\begin{equation}\label{eq:budget-attn-flops}
 F_{\mathrm{interaction}}\approx4BT^2h_qd_h=4BT^2d
\end{equation}
表示常规密集实现的两个主要矩阵乘法，不含投影与 Softmax。利用因果三角结构的专用实现可能减少实际算术，仅把未来分数设为负无穷的密集实现则照旧执行全部乘法。

令每层主干线性参数为 $P_\ell$，输入输出投影的执行参数按实际路径统计，整个前向可以近似写为
\begin{equation}
 F_{\mathrm{fwd}}\approx2BT\left(\sum_{\ell=1}^{L}P_\ell+Vd\right)
 +4LBT^2d+F_{\mathrm{elementwise}}.
\end{equation}
输入查表不按一次 $Vd$ 密集矩阵乘法计算，输出全词表投影却需要该项。权重共享减少存储，不删除输出计算。长上下文下第二项的占比显著上升，$2PBT$ 已吸收不了它。

以 $P_\ell\approx12d^2$ 的常见量级作粗略比较，注意力交互与该层线性前向之比约为 $\frac{T}{6d}$。长度相对宽度很大时，忽略交互项会系统性低估成本。具体转折位置随层结构和前馈宽度变化。

\subsection{重计算与有效工作量}
\term{激活重计算}{Activation Recomputation}{}通过丢弃部分中间量，在反向时重新执行相应前向，从而用计算换显存。若原训练前向计算为 $F_f$、反向为 $F_b$，额外重算为 $F_r$，实际计算为 $F_f+F_b+F_r$，此时全部成本记成固定 $6PD$ 会低估。

FlashAttention 等算法减少注意力中间矩阵的显式存储和数据搬运\cite{dao2022flashattention}，同时保留既定注意力的数学定义。使用同一公式计算的两个内核，可以因读写量、重算与硬件利用率不同而具有不同时间。资源报告应分开模型有效工作量、实际执行工作量与墙钟时间。

\section{推理显存与并发容量}

\subsection{权重存储与量化开销}
未量化权重存储为 $Pb_w$。若使用 $q$ 位量化且每 $g$ 个参数共享一个占 $b_s$ 字节的尺度，忽略边界取整时
\begin{equation}
 M_w\approx \frac{Pq}{8}+(\frac{P}{g})b_s+M_{\mathrm{other}}.
\end{equation}
$M_{\mathrm{other}}$ 包含零点、未量化层、对齐和打包元数据。某些后端还需临时反量化工作区。因此总显存由权重位宽与这些附加项共同决定，速度则取决于硬件支持、内核和数据移动。

\subsection{变长请求与键值缓存}
对同构 GQA，若第 $i$ 个并发请求已缓存 $T_i$ 个位置，则
\begin{equation}\label{eq:budget-kv}
 M_{\mathrm{KV}}=2Lh_{kv}d_hb_c\sum_{i=1}^{B}T_i.
\end{equation}
当全部请求长度相同才化为 $2LBTh_{kv}d_hb_c$。这里的长度包含已处理提示与已生成内容，只统计最终输出会大幅低估。共享前缀、滑窗淘汰、潜在状态和跨设备切分会改变实际驻留方式，需替换相应状态模型。

若分页块容量为 $p$ 个位置，且每条请求独立分配整块，则实际保留槽位为 $\sum_i p\lceil \frac{T_i}{p}\rceil$。内部浪费小于 $Bp$ 个位置，此外还存在页表与运行时保留。共享前缀时按逻辑引用数和物理页数各统计一次会造成重复计量。

\subsection{并发容量约束}
设设备可用显存为 $M_0$，权重及固定运行时占用 $M_f$，给定长度 $T$ 的单请求缓存为 $m(T)$，工作区上界为 $W(B,T)$，则并发必须满足
\begin{equation}
 M_f+B m(T)+W(B,T)\le M_0.
\end{equation}
若假设工作区固定，可先求缓存约束的整数上界；实际工作区随批处理变化时，需要在候选并发上逐项计算。平均请求长度描述典型状态，容量校验还需覆盖长请求同时到达的压力场景。

对于前述构造模型，取缓存每元素两字节，单请求每位置缓存为 $2(12)(4)(64)(2)=12288$ 字节。缓存8192个位置需要 $96$ MiB，32个同长请求需要 $3$ GiB。若键值头变成16，其缓存变成四倍；权重、工作区和服务时延仍需分别计算。

\section{训练状态、激活与分片}

\subsection{每参数字节数来自哪些状态}
一种混合精度 Adam 配置保存两字节训练权重、两字节梯度、四字节主权重以及两个四字节矩估计，总计每参数16字节。每参数开销由所保存的状态与精度共同决定。梯度累积可能使用四字节，某些实现不保留独立主权重，低精度优化器又会采用其他布局。

\begin{table}[htbp]
\centering
\caption{混合精度 Adam 训练中各状态分量每参数显存占用}
\label{tab:rev-training-state-memory}
\begin{tabular}{lcc}
\toprule 状态 & 本例每参数字节数 & 用途\\\midrule
计算权重 & 2 & 前向与反向的矩阵运算。\\
梯度 & 2 & 本次更新方向。\\
主权重 & 4 & 高精度参数累积。\\
一阶矩与二阶矩 & 8 & 自适应更新统计。\\\bottomrule
\end{tabular}
\end{table}
冻结基座不再为基座计算全部参数梯度和优化器状态，却通常仍需保存激活以向可训练适配器传递梯度。因此 LoRA 或 QLoRA 的显存无法由可训练参数百分比乘全参数训练显存得到。

\subsection{分片降低哪些驻留量}
设权重、梯度、优化器相关状态分别为 $M_w,M_g,M_o$，数据并行组有 $N_d$ 个设备。理想均匀分片的驻留模型为
\begin{align}
 M_{\mathrm{DDP}}&=M_w+M_g+M_o,\\
 M_{\mathrm{Z1}}&=M_w+M_g+\frac{M_o}{N_d},\\
 M_{\mathrm{Z2}}&=M_w+\frac{M_g+M_o}{N_d},\\
 M_{\mathrm{Z3}}&=\frac{M_w+M_g+M_o}{N_d}.
\end{align}
这里按 ZeRO 的分阶段思想区分被分片对象\cite{rajbhandari2020zero}；具体执行协议在\bookxref{ch:distributed-training}展开。公式没有包含临时聚合完整层参数的峰值、通信桶、预取和激活，也未保证所有层完全均匀。

若总参数状态112 GB，八设备理想全分片后平均为14 GB；某一执行阶段仍可能临时聚合一个大层。因此总状态除设备数给出的是驻留基线，单设备峰值还需在此基础上叠加临时项。

\subsection{激活模型的结构前提}
简单的线性保存模型可写为 $M_a\approx c_aLBTdb_a$，其中 $c_a$ 表示每层每位置保存多少个隐藏宽度等效元素。它应由实际计算图或测量校准，固定成一个普适常数会掩盖结构差异。若显式保存注意力概率，还可能增加 $O(LBh_qT^2)$ 元素；使用分块内核与重算后，保存结构发生变化。

峰值训练显存可组织为
\begin{equation}
 M_{\mathrm{train,peak}}=M_{\mathrm{resident}}+M_{\mathrm{activation}}
 +M_{\mathrm{gather}}+M_{\mathrm{communication}}+M_{\mathrm{workspace}}+M_{\mathrm{reserve}}.
\end{equation}
各项同时驻留的时间应按执行顺序分析。流水并行可能保存多个在途微批的激活；张量并行可能分摊部分张量，也可能复制其他状态。因此对全部分项统一除以设备总数会得到错误结果。

\section{带宽、时间与设备数}

\subsection{算术强度与不可突破的下界}
\term{算术强度}{Arithmetic Intensity}{}定义为操作数与相应存储层数据搬运字节数之比 $I=\frac{F}{Q}$。忽略重叠细节时，单个工作负载的耗时至少满足
\begin{equation}
 t\ge\max(\frac{F}{R_{\mathrm{peak}}},\frac{Q}{\beta}).
\end{equation}
一次调用若算术强度低，即使峰值矩阵吞吐很高，也可能受权重或缓存读取限制。提高批量可以复用权重，但同时增加缓存和延迟；峰值 FLOP/s 只是性能问题的一个维度。

\subsection{有效利用率与时间预算}
设整个训练所需模型工作量为 $F_{\mathrm{model}}$，$N$ 个设备的指定精度峰值总和为 $NR_{\mathrm{peak}}$。按模型工作量定义有效利用率
\begin{equation}
 u=\frac{F_{\mathrm{model}}}{N R_{\mathrm{peak}}t_{\mathrm{wall}}}.
\end{equation}
该定义将通信等待、数据等待、重算和其他开销反映在墙钟时间中。若另一个报告把重算计入分子，它采用不同口径，两个利用率无法直接比较。已包含通信损失的利用率再乘一次同样的通信折损系数，则会造成重复扣减。

在假设 $u$ 对候选设备数近似稳定时，期限 $S$ 要求
\begin{equation}
 N\ge\left\lceil\frac{F_{\mathrm{model}}}{R_{\mathrm{peak}}uS}\right\rceil.
\end{equation}
扩容后 $u$ 往往变化，因此严格预算需要 $u(N,\text{拓扑},B,T)$ 的情景模型。显存可行性也依赖具体并行方式，独立于计算约束另行确定。

\subsection{同一模型的训练与推理预算}
\begin{example}


\label{sec:budget-unified-example}
沿用前面 $L=12,d=1024,V=32000,h_q=16,h_{kv}=4,d_h=64,d_f=2816$ 的共享嵌入模型，总参数为 $P=168059904$。扣除归一化参数，参与主要线性映射的参数为 $P_{\mathrm{lin}}=168034304$。设处理约 $D=10^9$ 个非填充位置、序列长度 $T=2048$，忽略最后不满序列的边界修正。若每个独立块采用块内下一词元配对，有效监督量约为 $\frac{D(T-1)}{T}$；本段计算量使用处理位置数 $D$。不使用激活重算，注意力按稠密计算计数，则主要训练计算为
\begin{equation}
 F\approx6P_{\mathrm{lin}}D+12LTdD
 =1.310195712\times10^{18}\ \mathrm{FLOP}.
\end{equation}
第二项由每层每序列注意力前向 $4T^2d$ 乘以前反向近似因子3，再乘 $\frac{LD}{T}$ 得到。本例未计 softmax、逐元素运算和优化器更新，故是声明范围内的估算。

假设单设备峰值为 $2\times10^{13}$ FLOP/s，该计数口径下端到端有效利用率为0.3，则运行约 $218365.952$ 秒，即60.66小时。72小时期限下，计算量约束允许一个设备。容量还需独立检查：按每参数16字节保存训练状态，占2.5043 GiB；假设激活、工作区和余量合计4 GiB，峰值约6.5043 GiB，小于假定的16 GiB可用容量。

同一模型切换到推理时，两字节权重占320.55 MiB；四条请求各缓存2048个位置，KV占 $4\times2048\times12288=96$ MiB。若另预留512 MiB工作区，总计约928.55 MiB。推理吞吐还需结合推理阶段的形状、内核与带宽评估。由此，一组结构参数可以同时约束训练计算、训练驻留和推理缓存，但三个结果依赖不同的运行假设。
\end{example}

\subsection{资源联合预算}
\begin{example}


假设某稠密任务含 $P=7\times10^9$ 参数、训练 $D=10^{11}$ 词元，暂用 $F=6PD=4.2\times10^{21}$ 操作的基线。每设备峰值设为 $2\times10^{14}$ 操作每秒，模型利用率假设 $.4$，期限为10天，即864000秒。计算约束为
\begin{equation}
 N_{\mathrm{compute}}=\left\lceil
 \frac{4.2\times10^{21}}{(2\times10^{14})(.4)(864000)}
 \right\rceil=61.
\end{equation}
若拓扑要求设备数为八的倍数，可行候选从64开始。该设备数按上述工作量与利用率估算；长上下文交互和重算若未计入，还需增加相应工作量。

同一任务若采用本章16字节状态配置，总状态约104.31 GiB。设八路全分片，驻留状态约13.04 GiB，再假设激活、聚合、工作区和余量共8 GiB，则单设备约21.04 GiB，低于假设的40 GiB可用容量。该情景说明显存可能在八设备组内已经可行，而期限仍要求更大的集群。64设备的具体数据、张量和流水并行布局，以及其实际利用率，仍需由分布式执行方案确定。
\end{example}

\begin{algorithm}[构建资源可行情景]
\label{alg:budget-scenarios}
\AlgoInput{模型配置、词元预算、长度与并发分布、精度、候选拓扑、期限和容量余量。}
\AlgoOutput{每个候选的分项预算、可行条件和仍待测量的系数。}
\begin{enumerate}
\item 按实际共享关系计数参数，分别建立总参数、执行矩阵和特殊状态清单。
\item 对各拓扑计算权重、梯度、优化器、缓存、激活及峰值临时量；拒绝违反单设备容量约束的情景。
\item 按线性、注意力交互与重算分解工作量，使用该拓扑的利用率假设估计时间；检查期限约束。
\item 计算数据供给、检查点、CPU和主机内存需求，检查是否存在上游瓶颈。
\item 分别代入基准、保守和压力假设，记录结论改变的边界，而不止于输出一个设备数。
\end{enumerate}
测量系数存在不确定性时，可用区间构造预算情景。最终选型前用目标环境测量校准，方法见算力基础设施章。
\end{algorithm}

\section{主机、存储与数据供给预算}

训练吞吐目标若为每秒 $r_t$ 个词元，物化数据平均每词元占 $b_t$ 字节，则纯输入带宽下界为 $r_tb_t$；索引、压缩、重试和预取会增加其他开销。若运行时还需解析和分词，应使用对应步骤的实际处理速率估计 CPU 需求，按 GPU 数量乘固定核数得不到可靠结果。

检查点大小为 $C$，允许写入窗口为 $\Delta$，则平均有效写带宽至少为 $\frac{C}{\Delta}$。异步写入不减少数据，只是把等待转移到主机内存、队列和后台 I/O；产生速度长期大于写出速度时，队列仍会无限增长或触发阻塞。

总存储应分开原始数据、清洗数据、词元分片、检查点、评估结果、日志和中间资产。若保留 $K$ 个完整检查点，每个大小 $C$，至少需要 $KC$，另加构建暂存和恢复副本。分片文件数量影响元数据操作与小文件开销，容量充足不保证读取吞吐足够。

\begin{figure}[htbp]
\centering
\begin{tikzpicture}[>=Stealth,box/.style={draw,rounded corners,align=center,text width=28mm,minimum height=12mm,font=\small}]
\node[box] (data) at (0,0) {存储与数据队列\\字节/秒};
\node[box] (cpu) at (3.7,0) {CPU解析与预取\\样本/秒};
\node[box] (gpu) at (7.4,0) {设备计算与状态\\词元/秒、显存};
\node[box] (ckpt) at (7.4,-2) {检查点缓冲\\容量与写入窗口};
\draw[->] (data)--(cpu);\draw[->] (cpu)--(gpu);\draw[->] (gpu)--(ckpt);\draw[->] (ckpt.west)--(0,-2)--(data.south);
\end{tikzpicture}
\caption{资源供给与状态持久化形成闭环流水线}
\label{fig:rev-resource-data-path}
\end{figure}

\section{敏感性与预算表达}

预算报告应给出假设变化如何影响结论。例如固定模型时，普通KV缓存随总驻留词元线性变化；密集注意力交互随单序列长度平方变化；模型利用率从 $.4$ 降到 $.3$ 时，在其他条件不变的近似下，所需设备时间增加三分之一。参数量翻倍则同时影响状态与线性计算，但不必把所有缓存维度也翻倍。

余量要标明用途。容量安全余量、故障备用设备、维护窗口和负载峰值是不同项；若某系数已经按高分位负载测得，再把同一负载波动完整重复加入，会高估预算。反过来，为满足费用上限而把未知工作区当成零，预算就失去了保护作用。

最终预算是可追溯的情景集合：每个情景绑定模型版本、数据量、精度、长度分布、并行配置、有效速度与余量。它允许读者解释某项变化造成的后果，并指出哪些假设必须在实施时重新测量。




\chapterreferences
\myendchapterdocument
